\documentclass[conference]{IEEEtran}
\usepackage[T1]{fontenc}
\usepackage[utf8]{inputenc}
\usepackage{url}
\usepackage{amsmath}

\title{Single-Cycle RV32I Processor}

\author{
\IEEEauthorblockN{Viraj, Kaustubh, Pratyush, Bhargav and Rudraksh}
\IEEEauthorblockA{Department of Computer Applications / Computer Science and Engineering\\
Institution and Department: Confirm before formal submission}
}

\begin{document}
\maketitle

\begin{abstract}
This paper presents the design, microarchitectural implementation, rigorous multi-level verification, and open-source generic synthesis of a single-cycle processor executing a ten-instruction educational subset of the RISC-V RV32I base integer specification. Implemented in Verilog-2001, the processor integrates a 32-bit Program Counter, an asynchronous 256-word instruction read-only memory (ROM), an exact instruction decoder and control unit, a 32-word general-purpose register file with a hardwired-zero x0 register, a multi-format sign-extending immediate generator, an arithmetic and logic unit (ALU), a dedicated branch comparator unit, a 256-word data memory (RAM), and writeback multiplexing. Runtime defensive fault-gating logic suppresses state side-effects and latches diagnostic codes upon fetch misalignment, illegal instruction encodings, or invalid memory addresses. Functional correctness is established through 276,741 passing RTL assertion checks across 16 self-checking testbenches, supplemented by cycle-by-cycle comparison against an independent architectural reference model executing 38 instructions spanning all ten operations and four branch outcomes. Generic logic synthesis via Yosys yields 22,695 leaf generic cells with zero inferred latches. A port-observable post-synthesis netlist simulation passes 163 cycle-accurate checks, verifying hardware structure and confirming resolution of a synthesis ROM initialization defect. Results are strictly bounded to the tested instruction subset and generic cell library; no FPGA clock frequency, look-up table utilization, power consumption, or physical deployment is claimed.
\end{abstract}

\section{Introduction and Objectives}
The primary objective of this project is to provide a transparent, cycle-accurate reference implementation of a RISC-V processor where the fundamental stages of instruction execution---instruction fetch, instruction decode, operand selection, ALU execution, memory access, and register writeback---occur deterministically within a single clock cycle. By implementing a self-contained single-cycle microarchitecture, the processor eliminates the structural and data hazard complexities of pipelined designs while making datapath signal flow and control transitions directly inspectable in simulation waveforms.

The processor implements an educational subset of ten instructions from the RV32I base integer instruction set:
\begin{itemize}
\item \textbf{R-type (Arithmetic and Logic):} \texttt{ADD}, \texttt{SUB}, \texttt{AND}, \texttt{OR}, \texttt{XOR}
\item \textbf{I-type (Immediate Arithmetic and Load):} \texttt{ADDI}, \texttt{LW}
\item \textbf{S-type (Store):} \texttt{SW}
\item \textbf{B-type (Conditional Branch):} \texttt{BEQ}, \texttt{BNE}
\end{itemize}

To maintain an educational, fully verifiable scope, the microarchitecture deliberately omits other RV32I base instructions (shifts, comparison immediates, upper immediates, jumps, and byte/halfword memory accesses), privileged architecture (CSRs, privilege rings, architectural trap handlers), interrupts, pipeline stages, and cache hierarchies.

In lieu of complex operating system trap handlers, an educational fault-gating mechanism protects system state against invalid operations. Any misaligned instruction fetch, unrecognized opcode/funct field, out-of-bounds RAM access, or misaligned branch target immediately deasserts write enables (\texttt{reg\_write = 0}, \texttt{mem\_write = 0}), freezes the Program Counter, latches a 2-bit diagnostic fault code (\texttt{2'b01}: Fetch fault, \texttt{2'b10}: Illegal instruction fault, \texttt{2'b11}: Data/branch fault), and maintains this frozen state until cleared by an active-high synchronous reset.

\section{Literature Review}
The processor design and verification framework build upon primary specifications and established open-source Electronic Design Automation (EDA) tooling:

The RISC-V Instruction Set Manual, Volume I: Unprivileged ISA (Version 2.1, release v20260120) \cite{ref1} defines the programmer-visible contract: 32 general-purpose registers (\texttt{x0}--\texttt{x31}) where \texttt{x0} is hardwired to zero, fixed 32-bit instruction encodings, uniform register address placement (\texttt{rs1} at [19:15], \texttt{rs2} at [24:20], \texttt{rd} at [11:7]), signed two's-complement immediate layouts, and exact arithmetic/logical operation semantics. The specification decouples architectural definitions from microarchitectural realizations.

The Icarus Verilog simulation suite \cite{ref2} provides the event-driven simulation engine for functional verification. Testbenches apply stimulus vectors, inspect internal registers, evaluate bus assertions on clock edges, and export Value Change Dump (\texttt{.vcd}) waveforms for visual signal inspection.

The Yosys Open SYnthesis Suite \cite{ref3} provides RTL elaboration, generic technology mapping (\texttt{synth -top rv32i\_cpu}), and formal structural checking (\texttt{check -assert}). Synthesis verifies that behavioral Verilog constructs infer valid combinational gates and flip-flops without unwanted transparent latches or undriven nets.

Educational MIPS-32 processor architectures \cite{ref4} archived in project records served as structural inspiration for modular 5-stage partitioning (Fetch, Decode, Execute, Memory, Writeback). However, MIPS instruction encodings, control matrices, coprocessor registers, and simulation outputs are specific to MIPS and are not applicable to RISC-V. All RTL modules, decoders, immediate extractors, and testbenches in this repository are developed strictly for the RV32I specification.

\section{Proposed Solution and Implementation}
The single-cycle RV32I processor is organized into distinct, modular functional blocks interconnected by dedicated 32-bit buses.

A 32-bit Program Counter (\texttt{pc.v}) resets synchronously to \texttt{32'h00000000} and advances to \texttt{next\_pc} on each clock edge when unhalted. A 256-word instruction read-only memory (\texttt{instruction\_memory.v}) is combinationally addressed by \texttt{pc[9:2]}, returning the 32-bit instruction word within the cycle while enforcing 4-byte alignment (\texttt{pc[1:0] == 2'b00}) and range bounds (\texttt{0x000}--\texttt{0x3FC}).

The control unit and decoder (\texttt{control_unit.v}) decodes \texttt{opcode}, \texttt{funct3}, and \texttt{funct7}, generating write enables, multiplexer select lines, and ALU operation codes. Unrecognized instruction words deassert write enables and flag an illegal instruction fault. The register file (\texttt{register_file.v}) contains 32 registers with two asynchronous read ports and one synchronous write port; register \texttt{x0} discards writes and always reads zero. The immediate generator (\texttt{immediate_generator.v}) produces sign-extended 32-bit immediates for I-format, S-format, and B-format (with bit 0 implicitly zero).

The arithmetic logic unit (\texttt{alu.v}) performs addition, subtraction, AND, OR, and XOR, and computes effective memory addresses for \texttt{LW} and \texttt{SW}. A dedicated branch unit (\texttt{branch_unit.v}) compares register operands for \texttt{BEQ} and \texttt{BNE} and computes the branch target as \texttt{pc + imm\_b}. Data memory (\texttt{data_memory.v}) provides 256 words with asynchronous reads for single-cycle \texttt{LW} writeback and synchronous writes for \texttt{SW}. The writeback multiplexer (\texttt{writeback_mux.v}) selects between the ALU output and memory read data.

Top-level parameter \texttt{ARITHMETIC\_ONLY} is retained for backwards compatibility with Phase-3 regression tests. Programs terminate cleanly using a legal \texttt{BEQ} self-loop (\texttt{BEQ x0, x0, 0}).

\section{Verification Methodology}
The verification framework combines multi-tiered testbenches executed via \texttt{tools/run_tests.py}:
\begin{enumerate}
\item \textbf{Instruction Encoding Checks:} Validates machine-code word encodings against bitfield disassembly rules for all supported instructions.
\item \textbf{Independent Architectural Reference:} A Python reference emulator independently executes test programs, decodes instructions, updates state, and outputs golden cycle-by-cycle execution traces.
\item \textbf{Integration Test Suite (\texttt{tb\_full\_isa.v}):} Executes a 38-instruction program covering all ten supported operations, memory load/store sequences, and four branch outcomes (two taken, two not taken).
\item \textbf{Directed Component Testbenches:} Sixteen self-checking testbenches validate edge cases, including register file write priority, x0 protection on load, two's-complement arithmetic overflow, branch outcome truth tables, data memory alignment, synchronous reset during a pending store (\texttt{tb\_cpu_reset\_memory.v}), and fault injection.
\item \textbf{Runner Integrity Checks:} Test execution strictly validates output, rejecting runs with simulator errors, missing \texttt{PASS} banners, or timeouts.
\end{enumerate}

Key technical distinctions must be emphasized:
\begin{itemize}
\item \textbf{HDL Assertions vs. Executed Instructions:} The 276,741 reported assertion checks include 262,144 combinational checks from exhaustive sweeps of the instruction decoder, evaluating control lines across bitfield permutations. These represent combinational assertion evaluations, not 262,144 executed CPU instructions. The number of fully executed CPU instructions is bounded by the specific integration programs (e.g., 38 instructions in \texttt{tb\_full\_isa.v}).
\item \textbf{RTL Simulation vs. Netlist Simulation:} RTL simulation operates on behavioral Verilog with internal register visibility; netlist simulation operates on the gate-level netlist generated by Yosys synthesis and verifies only externally observable primary top-level ports (\texttt{pc}, \texttt{inst}, \texttt{alu\_result}, \texttt{mem\_write}, \texttt{reg\_write}, \texttt{fault}) over the 38-instruction trace.
\item \textbf{Bounded Testing vs. Formal Proof:} Bounded simulation proves correctness for the tested stimuli and edge cases; it does not constitute a mathematical formal equivalence proof across all possible 32-bit execution sequences.
\end{itemize}

\section{Results and Discussion}
All reported results map directly to verifiable test manifests and logs:
\begin{itemize}
\item \textbf{276,741 RTL assertion checks} passed across 16 test suites (\texttt{reports/verification/phase-5-manifest.json}).
\item \textbf{10 encoding vectors} verified (\texttt{reports/verification/phase-5-encoding\_checks.log}).
\item \textbf{38 reference instructions} executed covering 10 operations and 4 branch outcomes (\texttt{reports/verification/phase-5-phase4\_reference.log}).
\item \textbf{163 netlist checks} passed during port-observable post-synthesis simulation (\texttt{reports/synthesis/phase-5-netlist-check.json}).
\item \textbf{Structural synthesis checks} confirmed zero latches and zero residual memories/processes (\texttt{reports/synthesis/phase-5-stat.json}).
\end{itemize}

Generic logic synthesis using Yosys 0.71 produced:
\begin{itemize}
\item Hierarchical cell total: 22,704 (including 9 module instances).
\item Leaf generic logic cells: 22,695.
\item Generic flip-flop cells: 9,251 (Data RAM: 8,192; Register File: 1,024; PC/fault state: 35).
\item Multiplexer cells: 9,555.
\item Inferred latches: 0.
\end{itemize}

\subsection{Distinction Between Generic Cells and FPGA Resources}
Generic cells are abstract, technology-independent Boolean primitives. They do not correspond directly to FPGA Look-Up Tables (LUTs), Configurable Logic Blocks (CLBs), or dedicated hardware Block RAMs (BRAMs). The 8,192 flip-flops in data memory arise because single-cycle load writeback requires asynchronous read access; standard FPGA BRAMs require a synchronous read edge. Generic synthesis does not perform place-and-route, IO pad placement, clock tree synthesis, or device-specific timing closure.

\subsection{ROM Synthesis Defect and Regression}
During Phase-5 synthesis, an initialization flaw was uncovered: in behavioral simulation, ROM words were initialized to NOPs via a procedural loop before loading via \texttt{\$readmemh}. In Yosys structural synthesis, this procedural loop masked the \texttt{\$readmemh} contents, causing the synthesized ROM netlist to contain all NOPs. Netlist simulation detected this immediately on cycle 1, observing \texttt{00000013} instead of \texttt{04000093}. The lead corrected \texttt{instruction\_memory.v} by separating synthesis and simulation initialization paths. Following correction, the synthesized netlist passed all 163 port-observable checks over the 38-instruction trace. This RTL correction was authored by Viraj as part of the lead implementation; Rudraksh audited and verified the evidence.

\section{Performance Observations}
In this single-cycle processor, every valid non-faulting instruction executes and commits state updates on a single rising clock edge:
\begin{equation}
\text{CPI} = 1.0
\end{equation}
Fault holds and synchronous reset freeze execution rather than retiring instructions, and thus do not represent a multi-cycle CPI penalty.

A CPI of 1.0 does not imply high clock frequency. The achievable clock frequency ($F_{\max}$) is constrained by the long combinational propagation path. Testbench clock periods (e.g., 10~ns / 100~MHz) are arbitrary simulation timescale parameters and do not represent measured hardware frequency.

The critical path is the Load Word (\texttt{LW}) path:
\begin{equation}
T_{\text{clk}} \ge T_{\text{pc\_clk-to-q}} + T_{\text{ROM}} + T_{\text{decode}} + T_{\text{reg\_read}} + T_{\text{ALU}} + T_{\text{RAM}} + T_{\text{mux}} + T_{\text{setup}}
\end{equation}
Because instruction fetch, decode, register read, ALU effective address addition, RAM read, writeback multiplexing, and register setup all occur within one clock period, this long combinational delay inherently limits hardware clock speed. Determining an actual $F_{\max}$, setup/hold slack, or dynamic power requires physical synthesis and target-specific Static Timing Analysis (STA).

\section{Individual Contributions}
\begin{itemize}
\item \textbf{Viraj (Phase-1 \& Lead):} RTL implementation, test harnesses, simulation runners, Phase-1 fetch documentation, Phase-5 generic synthesis flow, and ROM initialization bug fix.
\item \textbf{Kaustubh (Phase-2):} Instruction-set documentation lead. Authored ISA reference documentation, instruction format field definitions, register operand usage tables, and worked address calculation examples in \texttt{docs/instruction-set/}.
\item \textbf{Pratyush (Phase-3):} Datapath and control documentation lead. Authored arithmetic CPU datapath documentation, control signal decoding matrices, reset/fault behavior, and ALU validation traces in \texttt{docs/architecture/}.
\item \textbf{Bhargav (Phase-4):} Verification and simulation documentation lead. Authored comprehensive test plan, simulation results analysis, waveform walkthroughs, test checklists, and Phase-4 handoff in \texttt{docs/verification/}.
\item \textbf{Rudraksh (Phase-5):} Academic reports and final submission artefacts lead. Audited project progress, reviewed and expanded Phase I report and final technical manuscript, verified primary references, authored demonstration script, and compiled interview preparation notes.
\end{itemize}

Initial Phase-5 documentation and LaTeX drafts were prepared with assistant support under the lead's identity. Rudraksh conducted an independent technical review, verified reported data against manifests and logs, synchronized manuscripts, and submitted these contributions under his verified GitHub identity (\texttt{Rudraksh <rudraksh969977@gmail.com>}).

\section{Conclusion}
The Single-Cycle RV32I Processor project demonstrates a modular, cycle-accurate implementation of a ten-instruction RISC-V subset in Verilog-2001. Multi-level verification—encompassing 276,741 passing RTL checks across 16 testbenches, 38 reference instructions executed cycle-by-cycle, and 163 passing port-observable netlist checks—confirms the functional correctness of the datapath, control unit, and defensive fault-gating logic. Generic Yosys synthesis validates structural integrity with 22,695 leaf generic cells and zero inferred latches. Remaining submission actions include team author sign-offs, recording the demonstration video, compiling the IEEE PDF on a TeX-enabled environment, and capturing graphical waveform screenshots if required.

\begin{thebibliography}{4}
\bibitem{ref1} RISC-V International, ``RV32I Base Integer Instruction Set, Version 2.1,'' in \emph{The RISC-V Instruction Set Manual, Volume I: Unprivileged ISA}, release v20260120, Jan. 2026. \url{https://docs.riscv.org/reference/isa/v20260120/unpriv/rv32.html}
\bibitem{ref2} S. Williams and contributors, ``Icarus Verilog Documentation and User Guide,'' Release Documentation, 2026. \url{https://steveicarus.github.io/iverilog/}
\bibitem{ref3} C. Wolf and YosysHQ, ``Yosys Command Reference and Synthesis Manual,'' YosysHQ Documentation, 2026. \url{https://yosyshq.readthedocs.io/projects/yosys/en/latest/using_yosys/synthesis/synth.html}
\bibitem{ref4} Project Team, ``Source Provenance and Architectural References,'' repository \texttt{REFERENCES.md}, Nov. 2026.
\end{thebibliography}

\end{document}
